\documentclass[11pt]{article}
\usepackage[T1]{fontenc}
\usepackage[utf8]{inputenc}
\usepackage{lmodern}
\usepackage[margin=1in]{geometry}
\usepackage{amsmath,amssymb,amsthm,mathtools}
\usepackage{booktabs,array,longtable}
\usepackage{microtype}
\usepackage{xcolor,listings}
\usepackage{xurl}
\usepackage[hidelinks]{hyperref}
\usepackage{enumitem}
\setlength{\emergencystretch}{2em}
\setlist{nosep}
\newtheorem{theorem}{Theorem}[section]
\newtheorem{lemma}[theorem]{Lemma}
\newtheorem{corollary}[theorem]{Corollary}
\newtheorem{proposition}[theorem]{Proposition}
\theoremstyle{definition}
\newtheorem{definition}[theorem]{Definition}
\theoremstyle{remark}
\newtheorem{remark}[theorem]{Remark}
\newcommand{\R}{\mathbb R}
\newcommand{\norminf}[1]{\left\|#1\right\|_{\infty}}
\newcommand{\code}[1]{\nolinkurl{#1}}
\newcommand{\leanref}[2]{\par\smallskip\noindent\textit{Lean correspondence.}
  \code{#1} in \code{#2}.\par\smallskip}
\lstset{basicstyle=\ttfamily\small,columns=fullflexible,
  keepspaces=true,breaklines=true,breakatwhitespace=false,
  showstringspaces=false,tabsize=4,aboveskip=0.7em,belowskip=0.7em}
\numberwithin{equation}{section}
\title{Interval Taylor Evaluation of \texorpdfstring{$yH[y_u]$}{yH[yu]}\\
  \large Mathematical estimates, Julia implementation, and Lean statements}
\author{}
\date{Repository snapshot: September 9, 2026}

\begin{document}
\maketitle
\begin{abstract}
This independent note explains how the term $P=yH[y_u]$ is evaluated in the
two-million-point Crunchy2 residual calculation. It develops the local
derivative enclosures, the seventh-order remainder bound, the role of bound
tightening, and the signed assembly of the residual. The relevant Lean
statements are presented as ordinary mathematical theorems with proofs.
The distinction between analytic theorems, hypotheses about numerical
enclosures, and the numerical assumptions of the dedicated certificate is
kept explicit. No part of the note requires another manuscript or external
LaTeX input file.
\end{abstract}

\section{The function and the certified inputs}
\label{sec:inputs}

Let $a_0,\ldots,a_d\in\R$ be the fixed coefficients of the candidate and set
\begin{equation}
 y(u)=\frac{a_0}{2}+\sum_{k=1}^{d}a_k\cos(ku),\qquad
 Q(u)=H[y_u](u),\qquad P(u)=y(u)Q(u).
 \label{eq:profile}
\end{equation}
The periodic Hilbert transform has the convention
\begin{equation}
 H[1]=0,\qquad H[\cos(ku)]=\sin(ku),\qquad
 H[\sin(ku)]=-\cos(ku)\quad(k\geq1).
 \label{eq:hilbert-convention}
\end{equation}
In particular $P=yH[y_u]$ is a pointwise product. It is different from the
other residual term $H[yy_u]$.

The run documented in
\code{wave-cap/BabenkoRigorous/results/phase4_scalar_n150_20260827_2m_20260829/README.md}
uses the following parameters:
\begin{center}
\begin{tabular}{lr}
\toprule
Quantity & Value\\
\midrule
Fourier degree $d$ & $150$\\
Number of coefficients & $151$\\
Number of grid points $n$ & $2{,}000{,}000$\\
Taylor order $M$ & $6$\\
Arb working precision & $128$ bits\\
Maximum rows per shard & $1024$\\
\bottomrule
\end{tabular}
\end{center}
Throughout, put
\begin{equation}
 h=\frac{2\pi}{n},\qquad u_j=(j-1)h\quad(1\leq j\leq n).
 \label{eq:grid}
\end{equation}
The cells $[u_j,u_j+h]$ cover $[0,2\pi]$, including its right endpoint.
The estimates use the full spacing $h$, not $h/2$.

The exact candidate is the one specified by the finite stored numerical
inputs, with their exact dyadic interpretation. The certificate is not a
claim about an unknown exact solution whose coefficients merely round to
these inputs. Arb supplies outward enclosures of all subsequent arithmetic,
including $\pi$, the grid locations, and trigonometric evaluations.

\begin{theorem}[Fourier identification and constant-mode convention]
\label{thm:fourier}
For the function in \eqref{eq:profile},
\begin{equation}
 y_u(u)=-\sum_{k=1}^{d}ka_k\sin(ku),\qquad
 Q(u)=\sum_{k=1}^{d}ka_k\cos(ku).
 \label{eq:fourier-Q}
\end{equation}
If $\alpha_0=a_0/2$ and $\alpha_k=a_k$ for $k\geq1$, the full-constant-mode
sum $\sum_{k=0}^{d}\alpha_k\cos(ku)$ equals $y$. Halving only the zeroth
coefficient changes neither the derivative sum nor its Hilbert transform.
\end{theorem}
\begin{proof}
Differentiate the finite sum term by term. Apply linearity of $H$ and
\eqref{eq:hilbert-convention} to each sine term to obtain
$H[-ka_k\sin(ku)]=ka_k\cos(ku)$. The coefficient conversion follows from
$\cos(0u)=1$. Both the derivative and the transformed derivative multiply
the zeroth coefficient by $0$, so neither depends on its value.
\end{proof}
\leanref{periodicHilbert_yu_eq_Hyu}{HilbertOnTrigPoly.lean}
The constant-mode statements correspond to \code{yHalfDC_apply},
\code{yu_halfDCCoeffs}, and \code{Hyu_halfDCCoeffs} in
\code{HalfDCConvention.lean}. Lean's \code{y}, \code{yu}, and \code{Hyu}
are defined in \code{Definitions.lean}; its coefficient vector has length
$d+1=151$. The generated half-DC certificate already makes the required
zeroth-coefficient conversion.

The product $P$ is a finite trigonometric polynomial of degree at most
$2d=300$, by the cosine product identity. The Julia path below forms its
derivative data without projecting it back to degree $150$. In particular,
this term requires neither an FFT nor an unresolved-mode estimate.

\section{Interval Taylor records and the input bounds}
\label{sec:records}

\begin{definition}[Valid Taylor record]
\label{def:record}
For $f\in C^{M+1}(\R)$, a valid Taylor record on a grid consists of a
positive spacing $h$, centers $c_{j,m}$, nonnegative errors $e_{j,m}$ for
$0\leq m\leq M$, and global bounds $B_m(f)$ for $0\leq m\leq M+1$, such that
\begin{align}
 |f^{(m)}(u_j)-c_{j,m}|&\leq e_{j,m}, &&0\leq m\leq M,\label{eq:record-grid}\\
 |f^{(m)}(u)|&\leq B_m(f), &&u\in\R,\quad0\leq m\leq M+1.
 \label{eq:record-global}
\end{align}
Write $I_{j,m}(f)=[c_{j,m}-e_{j,m},c_{j,m}+e_{j,m}]$ and
\begin{equation}
 U_{j,m}(f)=\sup\{|x|:x\in I_{j,m}(f)\}=|c_{j,m}|+e_{j,m}.
 \label{eq:U}
\end{equation}
\end{definition}

This is the mathematical content of \code{TaylorSeriesData} and
\code{IsTaylorSeriesOf} in \code{TaylorSeriesCorrectness.lean}. Lean uses
grid indices $0,\ldots,n-1$; Julia uses $1,\ldots,n$. Julia stores the whole
Arb ball in \code{derivs[j,m+1]}, while Lean's abstract record separates
its center and error. These are raw derivatives, not coefficients divided
by $m!$. There are seven local derivative columns for $M=6$, but eight
global bound entries, covering derivative orders $0,\ldots,7$.

An entry of Julia's \code{sup_bounds} is also an Arb value. It represents
a rigorously evaluated bound expression. Its upper endpoint is a usable
scalar upper bound; the ball is not asserted to enclose the exact supremum
norm. In the formulas below, scalar bounds describe the mathematical
inequalities. Julia implements the same bound expressions with outward
rounding; no claim of exact floating-point equality with an idealized
real-arithmetic formula is needed.

\begin{theorem}[Coefficient bounds and exact Taylor data]
\label{thm:coefficients}
For $f(u)=\sum_{k=0}^{d}(\alpha_k\cos(ku)+b_k\sin(ku))$, define
\begin{equation}
 C_m(f)=\sum_{k=0}^{d}k^m(|\alpha_k|+|b_k|),\qquad0^0:=1.
\end{equation}
Then $|f^{(m)}(u)|\leq C_m(f)$ for all $u\in\R$ and all $m\geq0$.
Consequently, a record with exact derivative samples, zero errors, these
global bounds, and positive spacing is a valid Taylor record.
For the factors $y,Q$ in \eqref{eq:profile}, valid bounds are
\begin{align}
 B_0(y)&=\frac{|a_0|}{2}+\sum_{k=1}^{d}|a_k|,&
 B_m(y)&=\sum_{k=1}^{d}k^m|a_k|\quad(m\geq1),\label{eq:y-bounds}\\
 B_m(Q)&=\sum_{k=1}^{d}k^{m+1}|a_k|&& (m\geq0).
 \label{eq:Q-bounds}
\end{align}
\end{theorem}
\begin{proof}
Each differentiation of a nonconstant sine or cosine multiplies its
coefficient by $k$ and changes only its sign and trigonometric type. Since
sine and cosine have absolute value at most one, the triangle inequality
gives the asserted sum. Constant terms vanish after one differentiation;
at order zero the convention $0^0=1$ supplies a valid bound, even for the
identically zero sine term. Finite trigonometric sums are smooth, so exact
samples with zero errors satisfy \eqref{eq:record-grid}, and the coefficient
estimate supplies \eqref{eq:record-global}. Substitute
$\alpha_0=a_0/2$, $\alpha_k=a_k$, and the coefficients of $Q$ to obtain
\eqref{eq:y-bounds}--\eqref{eq:Q-bounds}.
\end{proof}
\leanref{trigpoly_to_taylor_correct}{TaylorSeriesCorrectness.lean}
This Lean theorem assumes the identities giving the exact derivative
samples, zero sample errors, and the displayed coefficient-bound formula.
It is not a proof that the finite-precision Arb output has zero error.
In the numerical pipeline, inclusion of the exact derivative in the
computed ball is supplied by the interval evaluation instead.

Julia's \code{trigpoly_to_taylor} forms coefficient derivatives through
order $M+1$, evaluates orders $0,\ldots,M$ at the grid using \code{tp_eval},
and applies \code{supnorm_bound} to all orders $0,\ldots,M+1$.
Thus a seventh-order bound is available without sampling seventh
derivatives on the grid. The factors $y$ and $Q$ are enclosed separately
at each grid point, but no separate whole-cell supremum bound is taken
before their local derivative records are multiplied.

\section{Leibniz multiplication and its formal statements}
\label{sec:product}

\begin{theorem}[Leibniz derivative bound]
\label{thm:leibniz}
Suppose $f,g$ are $C^m$ near $u$, and
$|f^{(k)}(u)|\leq F_k$, $|g^{(k)}(u)|\leq G_k$ for $0\leq k\leq m$.
Then
\begin{equation}
 |(fg)^{(m)}(u)|\leq\sum_{k=0}^{m}\binom{m}{k}F_kG_{m-k}.
 \label{eq:leibniz-bound}
\end{equation}
If the factor bounds hold globally, the same expression is a global bound.
\end{theorem}
\begin{proof}
Repeated differentiation gives
$(fg)^{(m)}=\sum_{k=0}^{m}\binom{m}{k}f^{(k)}g^{(m-k)}$.
Take absolute values, apply the triangle inequality, and bound each factor.
The binomial coefficients and the factor bounds are nonnegative. The
argument holds at every point when the hypotheses hold globally.
\end{proof}
\leanref{abs_iteratedDeriv_mul_le}{LeibnizBound.lean}
The named Lean theorem is the pointwise statement; the global version is
its immediate application at every $u$.

For $P=yQ$, Julia constructs
\begin{align}
 I_{j,m}(P)&=\sum_{k=0}^{m}\binom{m}{k}
        I_{j,k}(y)I_{j,m-k}(Q), &&0\leq m\leq6,\label{eq:interval-product}\\
 L_m(P)&=\sum_{k=0}^{m}\binom{m}{k}B_k(y)B_{m-k}(Q),
        &&0\leq m\leq7.\label{eq:product-seed}
\end{align}
The interval operations in \eqref{eq:interval-product} round outwards.
For example, the order-two row uses the three terms
$I_{j,2}(y)I_{j,0}(Q)+2I_{j,1}(y)I_{j,1}(Q)+I_{j,0}(y)I_{j,2}(Q)$.
By interval inclusion and the Leibniz identity, this encloses $P''(u_j)$.
Equation~\eqref{eq:product-seed}, evaluated rigorously, supplies the initial
global bound vector. In particular its top entry includes all eight terms
of the seventh-order Leibniz sum.

\begin{theorem}[Correctness of a product Taylor record]
\label{thm:product-record}
Let valid records for $f,g\in C^{M+1}(\R)$ have the same grid, with centers
$c^f_{j,k},c^g_{j,k}$ and bounds $B_k(f),B_k(g)$. Define
\begin{equation}
 p_{j,m}=\sum_{k=0}^{m}\binom{m}{k}c^f_{j,k}c^g_{j,m-k}.
\end{equation}
Suppose numbers $\eta_{j,m}$ satisfy
$|(fg)^{(m)}(u_j)-p_{j,m}|\leq\eta_{j,m}$ for all
$1\leq j\leq n$ and $0\leq m\leq M$.
Then the centers $p_{j,m}$, errors $\eta_{j,m}$, and global bounds
$\sum_{k=0}^{m}\binom{m}{k}B_k(f)B_{m-k}(g)$ define a valid record for $fg$.
\end{theorem}
\begin{proof}
Products preserve $C^{M+1}$ regularity. The assumed error inequality is
exactly the sample-validity requirement. Theorem~\ref{thm:leibniz} supplies
the global bounds at each order. The common spacing remains positive.
\end{proof}
\leanref{mul_taylorseries_correct}{TaylorSeriesCorrectness.lean}
The inequality for $\eta$ is the explicit Lean hypothesis
\code{h_err_valid}; Lean does not infer the production Arb errors here.

To see how a sufficient error can be obtained, write exact input samples
as $c^f+\delta_f$ and $c^g+\delta_g$, with
$|\delta_f|\leq e^f$, $|\delta_g|\leq e^g$. Expanding one product gives
\begin{equation}
 |(c^f+\delta_f)(c^g+\delta_g)-c^fc^g|
 \leq |c^f|e^g+|c^g|e^f+e^fe^g.
\end{equation}
Summing these expressions with binomial weights gives a sufficient
$\eta_{j,m}$ for the exact center $p_{j,m}$. If the computed center is
$\widehat p_{j,m}$ and $|\widehat p_{j,m}-p_{j,m}|\leq\rho_{j,m}$,
add $\rho_{j,m}$ to the error. Julia's interval operations enclose both
input uncertainty and arithmetic rounding. This explanatory error formula
is not an assertion that Julia stores precisely that radius.

\begin{theorem}[Existence of exact records]
\label{thm:existence}
Every $2\pi$-periodic $C^{M+1}$ function has a valid exact-sample record
on any prescribed finite grid of positive spacing. Given valid records
for $f$ and $g$ on a common grid, there also exists a valid record for $fg$ with exact
samples, zero sample errors, and the Leibniz global bounds.
\end{theorem}
\begin{proof}
Each derivative through order $M+1$ of a periodic $C^{M+1}$ function is
continuous and periodic. Its absolute value attains a maximum on
$[0,2\pi]$, and periodicity extends this bound to $\R$. Take the true
derivatives as the centers and zero errors. For the second assertion,
take the true derivatives of $fg$ as centers and use
Theorem~\ref{thm:leibniz} for the supplied global bounds.
\end{proof}
\leanref{exists_isTaylorSeriesOf_of_contDiff_periodic}{CertificateAxioms.lean}
The specialized grid-spacing theorem is \code{cert_arb_trigpoly_to_taylor};
the product existence theorem is \code{cert_arb_ts_mul}. Despite their
names and their containing file, these declarations are theorems, not
numerical axioms. Their existential conclusions do not identify the witness
with the particular arrays produced by Crunchy2.

\section{From a Taylor record to a continuum estimate}
\label{sec:continuum}

\begin{theorem}[Taylor remainder on a cell]
\label{thm:cell}
Let $f\in C^{M+1}$ on a neighborhood of $[v,v+h]$, $h>0$, and suppose
$|f^{(M+1)}(u)|\leq B$ there. For $0\leq t\leq h$,
\begin{equation}
 \left|f(v+t)-\sum_{m=0}^{M}f^{(m)}(v)\frac{t^m}{m!}\right|
 \leq B\frac{t^{M+1}}{(M+1)!}.
 \label{eq:taylor-remainder}
\end{equation}
If $f^{(m)}(v)\in I_m$, then
\begin{equation}
 f(v+t)\in\sum_{m=0}^{M}I_m\frac{t^m}{m!}
 +\left[-B\frac{t^{M+1}}{(M+1)!},B\frac{t^{M+1}}{(M+1)!}\right].
\end{equation}
\end{theorem}
\begin{proof}
At $t=0$ the remainder is zero. For $t>0$, Taylor's theorem in Lagrange
form gives a point $\xi\in(v,v+t)$ for which the remainder is
$f^{(M+1)}(\xi)t^{M+1}/(M+1)!$. Apply the derivative bound. The interval
statement follows by enclosing each derivative value in $I_m$ and adding
the symmetric remainder interval.
\end{proof}
\leanref{abs_taylor_remainder_le}{TaylorRemainder.lean}
The Lean statement uses derivatives within the closed interval and a
Taylor polynomial defined with these derivatives. For the globally smooth
functions considered here they agree with the usual derivatives, giving
the displayed specialization. This avoids any endpoint regularity issue.

\begin{theorem}[Grid-to-continuum bound]
\label{thm:grid}
Let a valid Taylor record for $f$ be given on the grid \eqref{eq:grid}.
Then
\begin{equation}
 \sup_{u\in[0,2\pi]}|f(u)|\leq E(f):=
 \max_{1\leq j\leq n}\left\{
 \sum_{m=0}^{M}U_{j,m}(f)\frac{h^m}{m!}
 +B_{M+1}(f)\frac{h^{M+1}}{(M+1)!}\right\}.
 \label{eq:global-scan}
\end{equation}
For periodic $f$, the same bound holds on all of $\R$.
\end{theorem}
\begin{proof}
For any $u\in[0,2\pi]$, choose a cell with $u=u_j+t$ and $0\leq t\leq h$;
at $u=2\pi$ use the last cell. Apply Theorem~\ref{thm:cell}, the triangle
inequality, and $|f^{(m)}(u_j)|\leq U_{j,m}(f)$. Since $0\leq t\leq h$,
replacing $t^m$ by $h^m$ increases each nonnegative contribution. Taking
the maximum over $j$ proves the claim. Periodicity reduces any real
argument to a representative in one period.
\end{proof}
\leanref{grid_to_supNorm_taylor}{GridToSupNormTaylor.lean}
The record-level version is \code{sup_bound_taylorseries_correct} in
\code{TaylorSeriesCorrectness.lean}; it assumes that the displayed
cell expression is at most a supplied number for every $j$.

For $P$ and $M=6$, the local estimate is therefore
\begin{equation}
 |P(u_j+t)|\leq\sum_{m=0}^{6}U_{j,m}(P)\frac{h^m}{m!}
       +B_7(P)\frac{h^7}{7!},\qquad0\leq t\leq h.
 \label{eq:P-cell}
\end{equation}
This explains the continuum meaning of the product record. The production
calculation does not scan $P$ separately for the final residual certificate;
it uses the record in the signed assembly described in
Section~\ref{sec:residual}.

\section{What bound tightening does, and where it is used}
\label{sec:tightening}

\begin{lemma}[One grid-and-Lipschitz step]
\label{lem:lipschitz}
Let $f\in C^{m+1}(\R)$, suppose $|f^{(m+1)}|\leq C$ globally, and let
$A_m\geq\max_j|f^{(m)}(u_j)|$. Then
\begin{equation}
 |f^{(m)}(u)|\leq A_m+hC\qquad(u\in[0,2\pi]).
 \label{eq:lip-step}
\end{equation}
The same conclusion holds when the derivative bound is known only on the
covered interval. If $f$ is periodic, the result extends globally.
\end{lemma}
\begin{proof}
Choose $u=u_j+t$ with $0\leq t\leq h$. The mean value inequality for
$f^{(m)}$ gives $|f^{(m)}(u)-f^{(m)}(u_j)|\leq Ct$. Add the bound at $u_j$
and use $t\leq h$. Only values in the cell are used. For periodic $f$,
its derivatives are periodic, so one period suffices.
\end{proof}
\leanref{grid_lipschitz_sup_bound_step}{TightenSupBoundsLemma.lean}
The same file includes the local-interval variant used by the recursion.

\begin{theorem}[Downward grid-bound recursion]
\label{thm:recursion}
Suppose $f\in C^{M+1}(\R)$, $T_{M+1}\geq\norminf{f^{(M+1)}}$, and
$A_m\geq\max_j|f^{(m)}(u_j)|$ for $0\leq m\leq M$. Set, downwards,
\begin{equation}
 T_m=A_m+hT_{m+1},\qquad m=M,M-1,\ldots,0.
\end{equation}
Then $|f^{(m)}(u)|\leq T_m$ on the covered interval for every $m\leq M$.
If $f$ is periodic, the bounds are global.
\end{theorem}
\begin{proof}
The top-order hypothesis starts the induction. At order $m$, the bound
already proved at order $m+1$, together with Lemma~\ref{lem:lipschitz},
gives $A_m+hT_{m+1}$. Induct downwards and then use periodicity.
\end{proof}
\leanref{tighten_sup_bounds_correct}{TightenSupBoundsLemma.lean}
The recursion itself is \code{tightenSupBoundsRecursion}.

\begin{corollary}[The production update with a minimum]
\label{cor:production-tightening}
Let $L_m$ be valid global bounds and put
$A_m=\max_jU_{j,m}(f)$. For periodic $f$, the downward updates
\begin{equation}
 \widetilde B_{M+1}=L_{M+1},\qquad
 \widetilde B_m=\min\{L_m,A_m+h\widetilde B_{m+1}\}\quad(m=M,\ldots,0)
 \label{eq:production-tightening}
\end{equation}
are valid global bounds, with the top bound unchanged.
\end{corollary}
\begin{proof}
Induct downwards as in Theorem~\ref{thm:recursion}. Both entries of the
minimum are valid global bounds, so their minimum is also valid. Taking
outward upper endpoints preserves validity in interval arithmetic.
\end{proof}
This is a mathematical corollary describing Julia's
\code{_tighten_from_max!}; it is not being identified with a separate
Lean declaration of the same algorithm.

\begin{proposition}[The record transform currently defined in Lean]
\label{prop:lean-tightening}
Let a valid record for a periodic function have bounds $B_m$ and grid
absolute bounds $A_m=\max_j(|c_{j,m}|+e_{j,m})$. Define
\begin{equation}
 \begin{aligned}
 S_{M+1}&=B_{M+1},& S_m&=B_m+hS_{m+1},\\
 \widehat B_{M+1}&=B_{M+1},&
 \widehat B_m&=\min\{B_m,A_m+hS_{m+1}\}.
 \end{aligned}
 \label{eq:lean-tightening}
\end{equation}
Keeping the centers, errors, and spacing fixed, and replacing $B$ by
$\widehat B$, preserves validity of the Taylor record.
\end{proposition}
\begin{proof}
The old global bounds $B_m$ are also bounds at all grid points. Applying
Theorem~\ref{thm:recursion} with these as its grid bounds proves validity
of $S_m$. Lemma~\ref{lem:lipschitz} with the actual grid bound $A_m$ then
proves validity of $A_m+hS_{m+1}$. Take the minimum with the old valid
bound. Periodicity gives global validity; all other record conditions are
unchanged.
\end{proof}
\leanref{tighten_sup_bounds_taylorseries_correct}{TightenSupBoundsLemma.lean}
Equation~\eqref{eq:lean-tightening} transcribes the current definition of
\code{TaylorSeriesData.tighten} for a nonempty grid. Its internal recursion
uses the old global bounds. It is therefore not identical to Julia's
recursion \eqref{eq:production-tightening}, which reuses the updated
next-order bounds and the per-order grid maxima. Both constructions are
valid, but this Lean theorem is not an equality or execution-equivalence
proof for the Julia routine.

For $P=yQ$, the product routine calls tightening after constructing
\eqref{eq:interval-product} and \eqref{eq:product-seed}. Its dependency is
\begin{equation}
 I_{j,m}(P)\ \longrightarrow\ U_{j,m}(P)\ \longrightarrow\
 A_m(P)\ \longrightarrow\ \widetilde B_m(P).
\end{equation}
Tightening does not change $I_{j,m}$ or $U_{j,m}$. At $M=6$ it can improve
orders $0,\ldots,6$, but leaves $B_7(P)$ unchanged.

\begin{proposition}[No effect of this pass on the final residual scan]
\label{prop:no-effect}
Suppose $P$ is used only additively in the residual, and the final scan is
performed at order six. Replacing only $B_0(P),\ldots,B_6(P)$ by tighter
values, while preserving its local derivative balls and $B_7(P)$, does not
change the data consumed by that scan.
\end{proposition}
\begin{proof}
Signed addition and subtraction assemble local derivative balls from local
derivative balls, without using the global bound vector. The seventh-order
residual bound uses the seventh-order bounds of the summands. Both of these
inputs are preserved. Formula~\eqref{eq:global-scan} uses no lower-order
global bounds.
\end{proof}
This proposition is a dependency observation about the current Julia
pipeline, not a separately named Lean theorem. It explains why the generic
tightening call is redundant for the final use of this particular product.

It does not imply that grid tightening is unused in the full residual.
For example, with $D=(y_u^2+(1+Q)^2)^{1/2}$ and
$C_1=\partial_u(y_uD^{-1})$, the identity
$C_1^{(7)}=(y_uD^{-1})^{(8)}$ gives the valid bound
\begin{equation}
 \norminf{C_1^{(7)}}\leq B_7(C_1):=
 \sum_{k=0}^{8}\binom{8}{k}B_k(y_u)\widetilde B_{8-k}(D^{-1}).
 \label{eq:upstream-tightening}
\end{equation}
The eighth-order product bound supplies the seventh-order bound for
$C_1$. The reciprocal record
is constructed at a higher order, and its tightened bounds through order
eight feed this product. These improved intermediate bounds can therefore
improve $B_7(R)$, even though tightening does not appear as a separate term
in the final Taylor formula. The amount of improvement depends on which
candidate bounds are actually smaller.

\section{Signed residual assembly and the final estimate}
\label{sec:residual}

Write the residual as
\begin{align}
 R&=\mu Q-y-P-H[yy_u]+\beta C_1-\beta H[G_u],\label{eq:residual}\\
 D&=\sqrt{y_u^2+(1+Q)^2},\qquad
 G=(1+Q)/D,\qquad C_1=\partial_u(y_u/D).
\end{align}
The other branches supply their own valid Taylor records. Positivity of
$D$ is required for their reciprocal operations. This note analyzes the
$P$ branch; it does not establish those other records from scratch.

\begin{theorem}[Linear assembly of valid records]
\label{thm:assembly}
Suppose $f_1,\ldots,f_s$ have valid records on a common grid and
$R=\sum_{\ell=1}^{s}\lambda_\ell f_\ell$, with real constants
$\lambda_\ell$. Then a valid record for $R$ is given by
\begin{align}
 c^R_{j,m}&=\sum_\ell\lambda_\ell c^\ell_{j,m},&
 e^R_{j,m}&=\sum_\ell|\lambda_\ell|e^\ell_{j,m},\label{eq:assembly-local}\\
 B_m(R)&=\sum_\ell|\lambda_\ell|B_m(f_\ell).
 \label{eq:assembly-global}
\end{align}
Consequently the order-six residual estimate is
\begin{equation}
 \norminf{R}\leq
 \max_j\left\{\sum_{m=0}^{6}U_{j,m}(R)\frac{h^m}{m!}
                +B_7(R)\frac{h^7}{7!}\right\},
 \label{eq:final-residual}
\end{equation}
provided the residual is periodic.
\end{theorem}
\begin{proof}
Differentiation commutes with the finite linear combination. Subtract the
center in \eqref{eq:assembly-local}; the triangle inequality bounds the
remaining sum of sample errors by $e^R_{j,m}$. Applying the same argument
to the exact derivatives proves \eqref{eq:assembly-global}. Smoothness
and the common positive spacing are preserved. Apply
Theorem~\ref{thm:grid} to the resulting record.
\end{proof}
\leanref{scalar_mul_taylorseries_correct}{TaylorSeriesCorrectness.lean}
The same file proves \code{sub_taylorseries_correct} and
\code{sup_bound_taylorseries_correct}. Their structural composition for
$\mu Q-y-P-H[yy_u]$ is \code{path_a_residual_taylor_correct} in
\code{ResidualAssemblyCorrectness.lean}. Formula~\eqref{eq:assembly-local}
is in real arithmetic; the numerical code also encloses rounding of the
assembled centers and includes it in the output balls.

The order of operations matters for the size of the estimate. The verifier
subtracts $I_{j,m}(P)$ into the signed residual derivative interval before
taking its absolute bound. It does not add separately computed positive
supremum estimates for the nonlinear terms. Thus the centers can retain
cancellation among residual terms. The seventh-order global bound still
uses the triangle inequality. The product's tightened bounds of orders
zero through six play no role in either of these operations.

\section{The Julia code path}
\label{sec:julia}

All file paths in this section are relative to
\code{wave-cap/BabenkoRigorous/src/}. The excerpts use variables from the
existing driver; they are not separate executable scripts.

\subsection{Fourier construction and local records}
In \code{ResidualBound.jl}, \code{_residual_bound_full} constructs this branch:
\begin{lstlisting}
y   = cosine_series(y_coeffs; prec=prec)
yu  = deriv(y)
Hyu = hilbert(yu)

M_h = M_order + 2 + extra_orders
ts_y_h = trigpoly_to_taylor(y, N_grid, M_h; prec=prec)
ts_Hyu = trigpoly_to_taylor(Hyu, N_grid, M_order; prec=prec)
ts_y   = truncate_order(ts_y_h, M_order)

ts_y_Hyu = ts_y * ts_Hyu
\end{lstlisting}
The functions \code{cosine_series}, \code{deriv}, \code{hilbert},
\code{tp_eval}, and \code{supnorm_bound} are in \code{FourierSeries.jl}.
The initial high-order record for $y$ is shared with the $yy_u$ branch.
\code{truncate_order} selects the first seven derivative columns and the
first eight global bounds for this product; it does not approximate $y$
by a different trigonometric polynomial.

In \code{TaylorSeries.jl}, \code{Base.:*} computes
\eqref{eq:interval-product}, calls \code{_leibniz_sup_bounds} for
\eqref{eq:product-seed}, constructs the output record, and returns
\code{tighten_sup_bounds!(out)}. The latter consists of
\code{_grid_abs_max_ubs!}, which scans the stored intervals, followed by
\code{_tighten_from_max!}, whose loop runs from \code{M:-1:0}.
The entry \code{sup_bounds[M+2]}, representing derivative order $M+1$,
is untouched.

\subsection{The production sharded route}
The stage \code{_stage_p1!} in \code{Sharded.jl} uses the same row operations
on blocks of global grid indices:
\begin{lstlisting}
for si in shard_range
    jr = su.shards[si]
    ts_y_h = trigpoly_to_taylor(su.y, su.N_grid, su.M_h;
                               prec=p, j_range=jr)
    ts_Hyu = trigpoly_to_taylor(su.Hyu, su.N_grid, M_order;
                               prec=p, j_range=jr)
    ts_y = truncate_order(ts_y_h, M_order)
    ts_y_Hyu = ts_y * ts_Hyu
    _acc!(acc_yHyu, ts_y_Hyu.derivs, M_order)
    write_matrix_slices(store, "yHyu", si, ts_y_Hyu.derivs)
end
\end{lstlisting}
The global grid size, not the shard size, determines $u_j$ and $h$.
The accumulator keeps the largest upper endpoint of the absolute
derivative interval at each order. After merging worker accumulators,
\code{_stage_fft!} reconstructs the globally valid bound vector:
\begin{lstlisting}
sb_y_Hyu = _leibniz_sup_bounds(sb_y, sb_Hyu,
                              M_order, p, su.binom_M)
_tighten_from_max!(sb_y_Hyu, st.acc_yHyu.ubs,
                   M_order, su.h_arb, p)
\end{lstlisting}
Although this stage also performs FFT work for other terms, these two
operations for $yH[y_u]$ are scalar bound calculations. The tightening
uses maxima over the entire grid. Bounds obtained from a single shard's
maxima are not used as certificates over the whole period.

The stage \code{_stage_s3} reloads the stored rows, forms the residual with
the same signs as \eqref{eq:residual}, and applies \code{_supbound_scan}.
The code obtains its remainder correction from
\begin{lstlisting}
wR, remR = _taylor_step_weights(M_order, su.h_arb, p)
fboundR = Arblib.ubound(sb_R[M_order + 2] * remR)
\end{lstlisting}
Here \code{wR} contains outward evaluations of $h^m/m!$, and \code{remR}
contains $h^{M+1}/(M+1)!$. The scan evaluates the sum of absolute local
derivative balls, adds the remainder as a radius using \code{add_error!},
and compares upper endpoints across cells. The returned Arb ball's upper
endpoint is the certificate. Its midpoint alone is not the residual bound.

\subsection{Cost of this branch}
At order six the local product has
$\sum_{m=0}^{6}(m+1)=28$ Leibniz summands per row, or $56$ million over
the full grid. One reduction across the seven derivative columns visits
$14$ million entries. The scalar seed vector has $36$ Leibniz summands
over orders $0,\ldots,7$, and one downward scalar tightening pass has
seven steps. The sharded driver performs both local product reductions
and accumulator reductions; these counts do not claim that only one such
scan occurs in the whole program. The archived run does not provide a
separate timer for $yH[y_u]$, so no isolated wall time is inferred here.
The full product matrix need not be resident in memory at once.

\section{What the dedicated Lean certificate establishes}
\label{sec:trust}

The general analytic declarations described above have ordinary
mathematical proofs. Their hypotheses must still be connected to numerical
data. In particular, existence of a valid Taylor record does not establish
that a particular serialized Arb array is that record. The product
correctness theorem explicitly assumes its sample-error inequality, and
the grid-to-continuum theorem explicitly assumes its cellwise bound.

The dedicated certificate uses
\code{CertificatePhase4ScalarN150Phase420260827HalfDC.lean} for exact
candidate constants and endpoints. Its numerical interface is in
\code{CertificatePhase4ScalarN150Phase420260827HalfDCAxioms.lean}.
It assumes the following two statements, rather than proving them by
replaying the two-million-point computation in Lean:
\begin{align}
 d_*&>0,\qquad D(u)\geq d_*\quad(u\in\R),\label{eq:anchor-D}\\
 |R(u)|&\leq\varepsilon_*\quad(u\in[0,2\pi]),
 \quad\text{given \eqref{eq:anchor-D}.}\label{eq:anchor-R}
\end{align}
They are named \code{cert_arb_dmin_bound} and
\code{cert_arb_sup_bound_le_eps0}, respectively. The archived upper
endpoint is approximately $1.7792768524432498\times10^{-13}$, and the
denominator lower endpoint is approximately $0.34465883518530394$.
These are bounds for the full residual and denominator, not an isolated
numerical bound on $\norminf{P}$.

\begin{theorem}[Global consequence of the numerical certificate]
\label{thm:certificate}
Assume \eqref{eq:anchor-D}--\eqref{eq:anchor-R} for the certified candidate.
Then $D(u)\geq d_*$ and $|R(u)|\leq\varepsilon_*$ for every real $u$.
For the exact constants in the dedicated certificate, this implies
\begin{equation}
 |R(u)|\leq1.78\times10^{-13},\qquad D(u)\geq0.3446588
 \qquad(u\in\R).
\end{equation}
\end{theorem}
\begin{proof}
The denominator assertion is already part of \eqref{eq:anchor-D}.
The residual is $2\pi$-periodic: its factors, derivatives, and periodic
Hilbert transforms have that period. Choose a representative of $u$
modulo $2\pi$ in $[0,2\pi]$ and apply \eqref{eq:anchor-R} there.
For the decimal statements, compare the exact rational certificate
endpoints with $178/10^{15}$ and $3446588/10^7$, respectively.
\end{proof}
The periodic extension is \code{cert_residual_le} in the dedicated
axioms module. The decimal comparisons are proved by
\code{residual_bounded_by_phase4_n150_half_dc_decimal} and
\code{D_bounded_below_by_phase4_n150_half_dc_decimal} in
\code{MainPhase4ScalarN150Phase420260827HalfDC.lean}. The two numerical
assumptions remain assumptions of these conclusions; the displayed proof
does not replace their interval-arithmetic verification.

For comparison, the generic record construction in
\code{CertificateAxioms.lean} uses the product-existence theorem in the form
\begin{lstlisting}
cert_arb_ts_mul h_y h_Hyu (hh_y.trans hh_Hyu.symm)
\end{lstlisting}
where \code{h_y}, \code{h_Hyu} establish valid factor records and
\code{hh_y}, \code{hh_Hyu} identify their spacings with $2\pi/n$.
This generic construction is not the numerical dependency chain of the
dedicated half-DC theorem.

\section{Source correspondence}
\label{sec:sources}

The files listed below are part of the repository snapshot used for this
note. Julia paths are relative to \code{wave-cap/BabenkoRigorous/src/};
Lean paths are relative to \code{wave-cap/WaveCAP/WaveCAP/}.
The mathematical proofs in this note are readable presentations and stated
specializations of the Lean declarations, not new Lean source files.
Corollary~\ref{cor:production-tightening} and
Proposition~\ref{prop:no-effect} are explicitly identified implementation
consequences rather than claims of separately formalized declarations.

The referenced modules build with Lean \code{v4.29.0-rc6}. An axiom audit
of the listed analytic declarations reports only \code{propext},
\code{Classical.choice}, and \code{Quot.sound}. The dedicated numerical
conclusions additionally depend on the assumptions identified in
Section~\ref{sec:trust}.

\begingroup
\small
\begin{longtable}{@{}>{\raggedright\arraybackslash}p{0.27\textwidth}
  >{\raggedright\arraybackslash}p{0.68\textwidth}@{}}
\toprule
Purpose & File and principal declarations\\
\midrule
\endfirsthead
\toprule
Purpose & File and principal declarations\\
\midrule
\endhead
Fourier algebra & \code{FourierSeries.jl}: \code{cosine_series},
 \code{deriv}, \code{hilbert}, \code{tp_eval}, \code{supnorm_bound}.\\[0.6em]
Taylor arithmetic & \code{TaylorSeries.jl}: \code{trigpoly_to_taylor},
 \code{truncate_order}, \code{Base.:*}, \code{_leibniz_sup_bounds},
 \code{tighten_sup_bounds!}, \code{sup_bound}.\\[0.6em]
Residual implementation & \code{ResidualBound.jl} and \code{Sharded.jl}:
 \code{_residual_bound_full}, \code{_stage_p1!}, \code{_stage_fft!},
 \code{_stage_s3}.\\[0.6em]
Hilbert identity & \code{HilbertOnTrigPoly.lean}:
 \code{periodicHilbert_yu_eq_Hyu}.\\[0.6em]
Coefficient convention & \code{HalfDCConvention.lean}:
 \code{yHalfDC_apply}, \code{yu_halfDCCoeffs}, \code{Hyu_halfDCCoeffs}.\\[0.6em]
Product bound & \code{LeibnizBound.lean}:
 \code{abs_iteratedDeriv_mul_le}.\\[0.6em]
Record correctness & \code{TaylorSeriesCorrectness.lean}:
 \code{trigpoly_to_taylor_correct}, \code{mul_taylorseries_correct},
 \code{sup_bound_taylorseries_correct}.\\[0.6em]
Taylor theorem & \code{TaylorRemainder.lean}: \code{abs_taylor_remainder_le};
 \code{GridToSupNormTaylor.lean}: \code{grid_to_supNorm_taylor}.\\[0.6em]
Grid tightening & \code{TightenSupBoundsLemma.lean}:
 \code{grid_lipschitz_sup_bound_step}, \code{tighten_sup_bounds_correct},
 \code{tighten_sup_bounds_taylorseries_correct}.\\[0.6em]
Generic existence & \code{CertificateAxioms.lean}:
 \code{exists_isTaylorSeriesOf_of_contDiff_periodic},
 \code{cert_arb_trigpoly_to_taylor}, \code{cert_arb_ts_mul}.\\[0.6em]
Residual assembly & \code{ResidualAssemblyCorrectness.lean}:
 \code{path_a_residual_taylor_correct}.\\
\bottomrule
\end{longtable}
\endgroup

\end{document}
